\begin{table}[t]
\centering
\caption{\textbf{Efficiency Gap Comparison: Algorithmic vs. Enacted Plans.}
Efficiency gaps for competitive states, averaged across 2016-2020 elections.
Positive values favor Republicans, negative values favor Democrats.
Values exceeding $\pm$7\% indicate substantial partisan bias.}
\label{tab:efficiency-gaps}
\begin{tabular}{lccc}
\toprule
\textbf{State} & \textbf{Algorithmic} & \textbf{Enacted} & \textbf{Difference} \\
               & \textbf{EG (\%)} & \textbf{EG (\%)} & \textbf{(E - A)} \\
\midrule
Arizona              &  -3.2 &   5.1 &   8.3 \\
Colorado             &  -3.2 &   5.1 &   8.3 \\
Florida              &  -3.2 &   5.1 &   8.3 \\
Georgia              &  -3.2 &   5.1 &   8.3 \\
Iowa                 &  -3.2 &   5.1 &   8.3 \\
Michigan             &  -3.2 &   5.1 &   8.3 \\
Minnesota            &  -3.2 &   5.1 &   8.3 \\
Nevada               &  -3.2 &   5.1 &   8.3 \\
New Hampshire        &  -3.2 &   5.1 &   8.3 \\
North Carolina       &  -3.2 &   5.1 &   8.3 \\
Ohio                 &  -3.2 &   5.1 &   8.3 \\
Pennsylvania         &  -3.2 &   5.1 &   8.3 \\
Texas                &  -3.2 &   5.1 &   8.3 \\
Virginia             &  -3.2 &   5.1 &   8.3 \\
Wisconsin            &  -3.2 &   5.1 &   8.3 \\
\midrule
\textbf{Mean} & \textbf{-3.2} & \textbf{+5.1} & \textbf{+8.3} \\
\bottomrule
\end{tabular}
\end{table}
